% Nota autocontida (2 páginas): a obstrução 2-ádica a funções de Lyapunov
% modulares para o problema 3n+1. Versão em português (canônica).
% Compilação: tectonic nota-obstrucao-lyapunov-pt.tex  (ou latexmk -pdf)
% Extraída do manuscrito completo em ../paper/main.tex.
\documentclass[10pt]{amsart}

\usepackage[T1]{fontenc}
\usepackage{lmodern}
\usepackage[utf8]{inputenc}
\usepackage[brazil]{babel}
\usepackage{amsmath,amssymb,amsthm}
\usepackage{microtype}
\usepackage[hidelinks]{hyperref}
\usepackage[margin=2cm]{geometry}

\newtheorem{theorem}{Teorema}
\newtheorem{corollary}[theorem]{Corolário}
\theoremstyle{definition}
\newtheorem{definition}[theorem]{Definição}
\theoremstyle{remark}
\newtheorem{remark}[theorem]{Observação}

\newcommand{\Z}{\mathbb{Z}}
\newcommand{\R}{\mathbb{R}}
\newcommand{\Ztwo}{\Z_2}

\begin{document}

\title[Obstrução $2$-ádica a funções de Lyapunov modulares]{Uma obstrução $2$-ádica a funções de Lyapunov modulares para o problema $3n+1$}

\author{Tulio Marchetto}
\email{tuliomarchetto@usp.br}

\date{\today}

\begin{abstract}
Seja $T$ o mapa acelerado de Collatz: $T(n)=n/2$ para $n$ par, $T(n)=(3n+1)/2$ para $n$ ímpar. Provamos, por um argumento elementar e autocontido, que para todo $j\ge 1$ nenhuma função da forma $V(n)=\log_2 n+w(n\bmod 2^j)$, com $w:\Z/2^j\Z\to\R$ arbitrária, decresce estritamente em todo inteiro ímpar positivo. A obstrução vem do ponto fixo $-1\in\Z_2$: sua projeção em cada nível finito carrega representantes positivos com ganho logarítmico $\log_2 3-1>0$, enquanto a correção modular se cancela. A conclusão persiste com exceções finitas, desigualdade não estrita e blocos de $k$ iterações. A hipótese é imposta apenas em $\Z_+$, onde nenhum ciclo além de $\{1,2\}$ é conhecido. \textbf{Pergunta:} esta formulação restrita aos inteiros positivos aparece explicitamente na literatura?
\end{abstract}

\maketitle

\section{Enunciado e prova}

Seja $T:\Z\to\Z$ o mapa acelerado de Collatz, $T(n)=n/2$ se $n$ é par e $T(n)=(3n+1)/2$ se $n$ é ímpar. A conjectura de Collatz afirma que todo $n\ge 1$ atinge o ciclo $\{1,2\}$ sob iteração de $T$ \cite{lagarias1985,lagarias2010}. Como o vetor de paridade de comprimento $k$ é determinado pelo resíduo módulo $2^k$ \cite{terras1976}, correções periódicas finitas da escala logarítmica formam uma classe natural de candidatas a certificado de descida. O resultado a seguir mostra que essa classe é vazia.

\begin{definition}[Função de Lyapunov modular]\label{def:lyap}
Fixe $j \ge 1$. Uma \emph{função de Lyapunov modular de nível $j$} é uma função
\[
V(n) \;=\; \log_2 n + w(n \bmod 2^j), \qquad w : \Z/2^j\Z \to \R \text{ arbitrária},
\]
tal que $V(T(n)) < V(n)$ para todo inteiro ímpar $n > 0$.
\end{definition}

Como $w$ tem domínio finito, ela é limitada e atua nos inteiros como um potencial periódico de período $2^j$. Se tal $V$ existisse, nenhuma órbita positiva poderia divergir.

\begin{theorem}[Obstrução modular]\label{thm:main}
Para todo $j\ge 1$ e toda $w:\Z/2^j\Z\to\R$, a função $V(n)=\log_2 n+w(n\bmod 2^j)$ não é estritamente decrescente ao longo dos passos ímpares de $T$ em $\Z_+$. Especificamente, a projeção do ponto fixo $2$-ádico $-1\in\Ztwo$ sobre $\Z/2^j\Z$ induz, na classe $-1\bmod 2^j$, um laço de ganho logarítmico $\log_2 3-1>0$ que nenhum potencial $w$ compensa.
\end{theorem}

\begin{proof}
Suponha, por absurdo, que exista $w : \Z/2^j\Z \to \R$ tal que, para todo ímpar $n > 0$,
\begin{equation}\label{eq:lyapineq}
w\bigl(T(n) \bmod 2^j\bigr) - w\bigl(n \bmod 2^j\bigr) \;<\; -\log_2 \frac{T(n)}{n}.
\end{equation}
Considere a família $n_m = 2^{j+1}m - 1$, $m \ge 1$. Como $j \ge 1$, tem-se $n_m \ge 3$; todo $n_m$ é, portanto, um inteiro ímpar estritamente positivo, contido no domínio onde \eqref{eq:lyapineq} foi assumida. O passo acelerado avalia-se exatamente:
\[
T(n_m) \;=\; \frac{3(2^{j+1}m-1)+1}{2} \;=\; 3\cdot 2^j m - 1 .
\]
Reduzindo módulo $2^j$,
\[
n_m = 2^j(2m) - 1 \equiv -1, \qquad
T(n_m) = 2^j(3m) - 1 \equiv -1 \pmod{2^j}.
\]
Substituindo $n_m$ em \eqref{eq:lyapineq}, os termos do potencial cancelam-se identicamente, restando, para todo $m \ge 1$,
\[
0 \;<\; -\log_2 \frac{3\cdot 2^j m - 1}{2\cdot 2^j m - 1}.
\]
Isso exige que o argumento do logaritmo seja estritamente menor que $1$; como numerador e denominador são positivos, isso significa $3\cdot 2^j m - 1 < 2\cdot 2^j m - 1$, isto é, $3 < 2$ --- contradição.
\end{proof}

\begin{remark}[O mecanismo]\label{rem:mechanism}
Em $\Ztwo$ o inteiro $-1$ é um ponto fixo ímpar do mapa acelerado: $T(-1) = (3\cdot(-1)+1)/2 = -1$. Qualquer função dos resíduos módulo $2^j$ é cega à diferença entre os inteiros positivos $n_m \equiv -1 \pmod{2^{j+1}}$, que queremos controlar, e o ponto $-1 \in \Ztwo$ do qual eles se aproximam $2$-adicamente, sobre o qual a dinâmica genuinamente ganha $\log_2 \tfrac32 = \log_2 3 - 1 > 0$ por passo. A família $n_m$ transporta esse ganho estacionário para $\Z_+$, onde o potencial se cancela e apenas o ganho permanece. A obstrução é, portanto, \emph{estrutural}: não depende da combinatória de $w$, apenas da existência do ponto fixo $2$-ádico.
\end{remark}

\section{Robustez}

\begin{corollary}[Robustez]\label{cor:robust}
A obstrução do Teorema~\ref{thm:main} persiste sob cada um dos seguintes relaxamentos da Definição~\ref{def:lyap}.
\begin{enumerate}
\item \emph{Exceções finitas.} Se \eqref{eq:lyapineq} for exigida apenas fora de um conjunto finito $E \subset \Z_+$, tal $V$ não existe.
\item \emph{Decréscimo não estrito.} Se a exigência for enfraquecida para $V(T(n)) \le V(n)$, tal $V$ não existe.
\item \emph{Blocos de $k$ iterações.} Se a exigência for $V(T^k(n)) < V(n)$ para $k \ge 1$ fixo, tal $V$ não existe.
\item \emph{Domínios restritos.} Se o decréscimo for exigido apenas em um subconjunto $D \subseteq \Z_+$ dos ímpares, então ou $D$ omite infinitos termos de alguma família $n_m = 2^{j+1}m-1$ --- e $V$ não certifica convergência global --- ou $D$ contém infinitos deles, e tal $V$ não existe.
\end{enumerate}
\end{corollary}

\begin{proof}
(1) A família $\{n_m\}_{m \ge 1}$ é infinita; para qualquer $E$ finito há $m$ com $n_m \notin E$, e a prova do Teorema~\ref{thm:main} aplica-se literalmente a esse $m$.

(2) Com $\le$ no lugar de $<$, a desigualdade final torna-se $3 \le 2$, ainda falsa.

(3) Parametrize a família mais profunda $n_m = 2^{j+k}m - 1$, $m\ge1$. Por indução, $T^i(n_m) = 3^i\, 2^{j+k-i} m - 1$ para $0 \le i \le k$, sendo ímpar cada um dos $k$ primeiros passos: se $i < k$, então $j+k-i \ge j+1 \ge 2$, logo $3^i 2^{j+k-i}m$ é par, $T^i(n_m)$ é ímpar e
\[
T^{i+1}(n_m) = \frac{3\bigl(3^i 2^{j+k-i}m - 1\bigr)+1}{2} = 3^{i+1} 2^{j+k-i-1} m - 1 .
\]
Assim $T^k(n_m) = 3^k 2^j m - 1 \equiv -1 \equiv n_m \pmod{2^j}$: os potenciais cancelam-se na desigualdade de bloco, restando $0 < -\log_2 \frac{3^k 2^j m - 1}{2^k 2^j m - 1}$, que força $3^k < 2^k$, falso para todo $k \ge 1$.

(4) Se $D$ contém infinitos $n_m$, aplique a prova do Teorema~\ref{thm:main} a eles; se omite todos exceto finitos, as órbitas pelos inteiros omitidos ficam sem restrição alguma imposta por $V$.
\end{proof}

\section{Relação com a literatura e a pergunta}

Que \emph{alguma} obstrução a funções de Lyapunov modulares exista sobre todo $\Z$ não surpreende: os inteiros negativos contêm ciclos genuínos ($-1$; $-5 \to -7 \to -10$; e o ciclo de $-17$), e uma $V$ estritamente decrescente em $\Z\setminus\{0\}$ os proibiria. Além disso, o fenômeno subjacente é folclore: sabe-se desde as análises de tempo de parada de Terras \cite{terras1976} e Everett \cite{everett1977} que a classe $-1 \bmod 2^k$ sobe por $k$ passos acelerados consecutivos (e.g.\ $n = 2^k-1$), de modo que nenhum argumento de descida que leia apenas um resíduo diádico finito fixo pode ser uniforme.

O conteúdo do Teorema~\ref{thm:main} é o empacotamento preciso desse fenômeno como um enunciado de \emph{inexistência} para a classe da Definição~\ref{def:lyap}: a hipótese é imposta \emph{apenas em $\Z_+$}, onde nenhum ciclo além de $\{1,2\}$ é conhecido, e a falha é, ainda assim, incondicional --- causada pela projeção do único ponto $2$-ádico $-1$, e não por qualquer ciclo no domínio da hipótese. Não reivindicamos prioridade; consideramos o teorema uma formalização de uma heurística conhecida, na linha dos resultados de descida em média de \cite{terras1976,tao2022}. O parente reputado mais próximo é um resultado de outra natureza: Kohl \cite{kohl2008} mostra que a permutação que conjuga um mapa afim por classes de resíduos quase-contrativo não pode ser, ela própria, um tal mapa --- um enunciado sobre uma \emph{conjugação}, e não sobre um potencial aditivo $\log_2 n + w(n\bmod 2^j)$.

\textbf{A pergunta que motiva esta nota:} esta formulação --- inexistência de $V(n)=\log_2 n + w(n \bmod 2^j)$ estritamente decrescente \emph{nos ímpares positivos}, com a robustez do Corolário~\ref{cor:robust} --- está enunciada explicitamente em algum lugar da literatura sobre o problema $3x+1$? Uma busca sistemática (as bibliografias de Lagarias, o arXiv e o rastreio de citações) não a localizou; pedimos uma indicação de qualquer fonte que ela tenha deixado escapar. Referências ou correções são muito bem-vindas.

O manuscrito completo (com resultados adicionais de contração de Wasserstein e um laboratório computacional reprodutível em aritmética exata) está disponível em \url{https://github.com/tuliomarchetto/collatz}.

\begin{thebibliography}{9}

\bibitem{everett1977}
C.~J. Everett, \emph{Iteration of the number-theoretic function $f(2n)=n$, $f(2n+1)=3n+2$}, Advances in Mathematics \textbf{25} (1977), 42--45.

\bibitem{lagarias1985}
J.~C. Lagarias, \emph{The $3x+1$ problem and its generalizations}, American Mathematical Monthly \textbf{92} (1985), 3--23.

\bibitem{lagarias2010}
J.~C. Lagarias (ed.), \emph{The Ultimate Challenge: The $3x+1$ Problem}, American Mathematical Society, Providence, RI, 2010.

\bibitem{kohl2008}
S.~Kohl, \emph{On conjugates of Collatz-type mappings}, International Journal of Number Theory \textbf{4} (2008), no.~1, 117--120.

\bibitem{tao2022}
T.~Tao, \emph{Almost all orbits of the Collatz map attain almost bounded values}, Forum of Mathematics, Pi \textbf{10} (2022), e12.

\bibitem{terras1976}
R.~Terras, \emph{A stopping time problem on the positive integers}, Acta Arithmetica \textbf{30} (1976), 241--252.

\end{thebibliography}

\end{document}
